\section{Research Objectives} \label{section:research_objectives}

\label{para:objectives_mli_goal} % ---------- %

This PhD aims to develop a machine-learned interface (MLI) model that, given two input materials; whether slabs, bulks,
or combinations, can generate realistic and physically plausible interface structures between them\nocite{}.

The purpose of this model is not to directly predict interfacial energies or properties, but to output candidate
geometries that are likely to be realisable and energetically reasonable, enabling downstream evaluation via existing
machine-learned potentials (MLPs) or density functional theory (DFT)\nocite{}.

\label{para:objectives_mli_pipeline} % ---------- %

The ultimate goal is to accelerate interface discovery by replacing brute-force combinatorial stacking with an
intelligent generative model\nocite{}.

By learning patterns in crystallography, registry, and bonding across known interfaces, the MLI model will propose
viable configurations that reduce the need for exhaustive enumeration and manual filtering\nocite{}.

\label{para:objectives_mli_extended_use} % ---------- %

Such a tool could eventually support applications beyond interface prediction, including surface termination prediction
under reactive gas atmospheres, simulation of crystal growth through layer-by-layer stacking, and real-time assessment
of growth defects during slab assembly\nocite{}.

\label{para:objectives_dataset_building} % ---------- %

To support development of the MLI model, the first phase of this research focuses on constructing a large, diverse
dataset of heterostructure interfaces\nocite{}.

Since direct generation and relaxation of thousands of interfaces using DFT would be prohibitively expensive,
machine-learned interatomic potentials (MLPs), specifically the MACE framework, are evaluated for their ability to
replicate DFT-derived rankings of interface favourability\cite{Batatia2022-xz}.

\label{para:objectives_screening_workflow} % ---------- %

If MLPs can reliably identify promising structures, even without full relaxation, they can serve as a low-cost screening
tool to filter candidate interfaces for inclusion in the training set\nocite{}.

This enables the creation of a scalable data pipeline where realistic interfaces are generated via tools like ARTEMIS,
screened using MLPs, and selectively validated with DFT\cite{Taylor2020-nl}.

\label{para:objectives_dft_validation} % ---------- %

Once promising interfaces are identified, a subset will be fully relaxed using DFT to provide high-fidelity data for MLI
training\nocite{}.

These DFT-relaxed structures will serve as the reference for extracting physically meaningful descriptors, such as
surface orientation, lattice mismatch, interfacial strain, atomic registry, and coordination environments, that can be
used as features for the MLI model\cite{Gouveia2025-fpml}.

\label{para:objectives_dft_mlp_comparison} % ---------- %

Although initial results show strong rank correlation between MLP and DFT energies (e.g., $\rho \approx 0.98$ for Sn|Ge), this
agreement may degrade in more complex systems or larger interface models.

Full DFT validation remains essential to calibrate and benchmark the workflow, particularly for low-symmetry or poorly
coordinated interfaces where MLP generalisation may falter.

\label{para:objectives_relaxation_scaling} % ---------- %

To further extend this validation, rankings based on MLP-relaxed geometries will be compared to those from DFT-relaxed
structures\nocite{}.

If similar rank orderings are recovered across both relaxation methods, it would support the use of MLPs for handling
much larger and more realistic interfacial systems, ones that are computationally infeasible to relax using DFT alone\nocite{}.

This opens the possibility of generating structurally diverse, high-volume datasets beyond the reach of traditional ab
initio approaches\nocite{}.

\label{para:objectives_inverse_design} % ---------- %

A longer-term extension of this research, contingent on the success of the main MLI model and access to further
computational resources, is the development of an inverse interface design framework\nocite{}.

Rather than proposing interfaces from input slabs alone, this model would take target properties or desired interfacial
features, such as low strain, specific registry types, or stability thresholds, and suggest plausible A|B configurations
that meet those criteria\nocite{}.

\label{para:objectives_inverse_limitations} % ---------- %

Such a tool could support experimental design by guiding researchers toward synthesizable interfaces that meet
functional requirements\nocite{}.

However, training such a model would require a much broader dataset, including properties computed from full DFT
relaxations on MLI-generated structures, something only feasible with additional funding and computational time\nocite{}.

\label{para:objectives_mlp_refinement} % ---------- %

To improve the accuracy of interface energy predictions, particularly for systems where MACE performs poorly (e.g.,
strained silicon-rich structures or cross-material environments), this project will explore refinement strategies for
the MLP itself\nocite{}.

One approach is delta-learning: training a secondary model to correct systematic energy differences between MACE and
DFT, using DFT-relaxed structures as a reference. \cite{Pitfield2025-ah}

\label{para:objectives_mlp_tuning} % ---------- %

Another strategy is direct fine-tuning or retraining of MACE on an expanded dataset that includes interface-specific
configurations, such as strained slabs, low-symmetry orientations, or cross-material bonding.

This is motivated by known failure modes in MLIPs, including "PES softening", poor generalisation to novel bonding
environments, and underprediction of surface and interface energies. \cite{}

These efforts will aim to mitigate known risks; such as extrapolation errors, by enriching the training distribution
with interfacial scenarios that challenge model stability.

By explicitly introducing diverse bonding motifs and strain conditions, the model's predictive performance can be made
more robust across real-world interface configurations, reducing the likelihood of catastrophic errors in
high-throughput screening\nocite{}.

\label{para:objectives_failure_modes} % ---------- %

A further objective of this research is to identify and mitigate failure modes in machine-learned potentials when
applied to interfacial systems\nocite{}.

Despite strong average performance, modern MLIPs such as MACE exhibit known limitations, including potential energy
surface (PES) softening, poor extrapolation to strained or defective configurations, and underprediction of surface
energies\cite{}.

\label{para:objectives_failure_diagnostics} % ---------- %

This work will systematically assess when and why these failures occur, starting with Group IV interfaces as test cases.

Specific attention will be paid to bonding asymmetries, coordination anomalies, and cases where poor generalisation
leads to incorrect rank ordering or unrealistic relaxations.

Where needed, refinement strategies such as delta-learning, fine-tuning, and active learning will be deployed to correct
these behaviours\cite{}.

The goal is not only to improve model accuracy, but also to build a deeper understanding of where and why MLPs break
down when applied to out-of-distribution geometries; informing the development of more robust, uncertainty-aware
workflows for interface prediction\nocite{}.

\label{para:objectives_structure_insight} % ---------- %

In parallel with predictive modelling, the project seeks to uncover the structural and chemical features that determine
interfacial stability\nocite{}.

One key line of enquiry is whether surface favourability correlates with interface favourability; for example, whether
low-energy surfaces tend to form stable interfaces, or whether compensating effects between unfavourable surfaces can
yield unexpectedly favourable interfaces\nocite{}.

\label{para:objectives_structure_analysis} % ---------- %

Further investigations will explore how stability varies across crystal types (e.g. fcc, hexagonal, tetragonal),
material classes (e.g. semiconductors, metals, insulators), and bonding mechanisms (e.g. covalent, ionic, metallic)\nocite{}.

The influence of valence orbitals; such as $s$ versus $p$ dominance, will also be considered, particularly in how they
affect registry, relaxation, and bonding across the interface\nocite{}.

\label{para:objectives_summary} % ---------- %

Together, these objectives form a cohesive research trajectory: from benchmarking machine-learned potentials and
constructing curated datasets, to training a generative interface model (MLI) and exploring its extension toward inverse
design\nocite{}.

The project aims not only to reduce the computational cost of interface exploration, but also to uncover structural
principles and predictive strategies that generalise across materials systems\nocite{}.

\label{para:objectives_impact} % ---------- %

By integrating machine learning with crystallographic modelling and ab initio validation, this work supports the
development of practical tools for scalable, data-driven interface discovery in electronic, optoelectronic, and
catalytic materials\nocite{}.

% ---------- %
